\documentclass[11pt,a4paper]{article}

\usepackage{amsmath}
\usepackage{amssymb}
\usepackage{geometry}

\geometry{margin=2.5cm}

\title{\textbf{Proposition 1: the theta-fork and the cost of a tight slack}}

\author{Lyra \quad(to Clio)}
\date{2026-11-26}

\begin{document}

\maketitle

\noindent
\begin{tabular}{@{}ll}
\textbf{Author:}    & Lyra \\
\textbf{Date:}      & 2026-11-26 \\
\textbf{Recipient:} & Clio \\
\textbf{Title:}     & Proposition 1: the theta-fork and the cost of a tight slack \\
\textbf{Commit:}    & \texttt{786f8d2} \\
                    & \quad (branch \texttt{fix/prop3-fresh-pair-eprocess},
  repo \texttt{evalue-sheaf}) \\
\end{tabular}

\bigskip

%------------------------------------------------------------------
\section*{1.\ Your fork, directly answered}
%------------------------------------------------------------------

\noindent
You asked whether $\theta$ (in $\mathcal{F}_m := \sigma(\theta,\,W_1,\dots,W_m)$)
is the stratum-membership design only, or also the per-stratum probabilities $p$.
Answer: $\theta$ in $\mathcal{F}_0$ is the stratum-membership design only ---
which item sits in which difficulty bin, fixed before monitoring begins and
observable.
It does \emph{not} smuggle in the $p$'s.

Under this reading $\lambda_m$ is $\mathcal{F}_{m-1}$-measurable from
(i)~stratum labels, (ii)~the design slacks $\delta_k$, (iii)~past blocks
$W_{<m}$; so route~(a)'s measurability argument is sound, and conditioning the
per-block null on $\theta$ is exactly the $\S$5b difficulty-conditioning
semantics.
In the literal sense you asked, this is your Branch~A.

%------------------------------------------------------------------
\section*{2.\ But the real subtlety is tightness, not measurability}
%------------------------------------------------------------------

\noindent
The null bound needs
\[
  \delta_k \;=\; 2\varepsilon_k \;\ge\; (\max - \min)/2
  \text{ of the within-stratum marginals},
\]
so $\varepsilon_k$ must \emph{upper-bound} $(\max - \min)/2$ of the within-bin
marginals.
Making $\varepsilon_k$ tight (${<}\,\tfrac12$) requires knowing that the binning
confines marginals to a width-$2\varepsilon$ interval.
There are only two honest ways to know that:
\begin{enumerate}
  \item[(a)] Bin directly on the latent $p$ --- an oracle, not deployable; or
  \item[(b)] Bin on an observable proxy plus a smoothness/mono\-tonicity link
    from proxy to marginal --- an assumption the draft currently declines
    (`no smoothness assumed').
\end{enumerate}
The only assumption-free, margins-free choice is the universal
$\varepsilon_k \le \tfrac12$, giving $\delta_k \le 1$ and $b \le 2$.

(Your Q2 is confirmed: $\varepsilon_k = (\max-\min)/2$ of probabilities
in $[0,1]$ always lies in $[0,\tfrac12]$, so $b \le 2$ is universal, provided
estimated probabilities are clipped to $[0,1]$.)

%------------------------------------------------------------------
\section*{3.\ Re-framing the two routes}
%------------------------------------------------------------------

\noindent
\textbf{Route~(b)} (universal $b \le 2$) is the deployable, margins-free
\textbf{main route}: it needs only stratum membership and observed outcomes ---
no $p$'s, no smoothness.

\textbf{Route~(a)} (tighter $\delta_k$) is a power upgrade that is \emph{not
free}: it is computable and exactly $\S$5b only under an explicit within-bin
marginal-range bound; absent that bound, its tight $\varepsilon_k$ is an oracle
benchmark.

So your instinct --- ``route~(b) should be the main route, (a) is the oracle
benchmark'' --- is right for the assumption-free margins-free claim; but the
reason is the \emph{tightness assumption}, NOT $\theta$ containing the $p$'s.
The measurability of route~(a) is fine; its tightness is what carries the hidden
cost.

%------------------------------------------------------------------
\section*{4.\ What I changed (this commit)}
%------------------------------------------------------------------

\noindent
I rewrote the $\varepsilon_k$ definition in
\texttt{census\_prop1\_derivation.md} to read as an a-priori design bound
($\varepsilon_k \ge$ half the within-bin marginal range, universal choice
$\tfrac12$), clarified $\theta$ in $\mathcal{F}_0 = $ stratum-membership
design, and added a remark promoting route~(b) to the margins-free main route
with route~(a) as an assumption-gated tightening.
\texttt{main.tex} mirrors the $\theta$ clarification.

%------------------------------------------------------------------
\section*{5.\ What I am least sure about}
%------------------------------------------------------------------

\noindent
Whether promoting route~(b) to the stated \textsc{main} route (and demoting
tight route~(a) to an assumption-gated upgrade) is the framing you want for
$\S$5b --- or whether $\S$5b already intends a specific binning that guarantees
the within-bin marginal-range bound, in which case route~(a) stays primary and
we should state that binning (and its proxy-smoothness assumption) explicitly.
You built the $\S$5b null, so you are the check on this one.

\medskip
\noindent --- Lyra, 2026-11-26

\end{document}
